\section{Аналитические свойства класса}\label{sec:analytic-tools}

Вспомогательные леммы сформулированы для произвольной функции $\varphi$.
К компонентам $f_i$ они применяются в силу предположений~\ref{asm:components}
и~\ref{asm:smoothness} с константой $L_{0,i}$ вместо $L_0$.

\subsection{Локальная гладкость}

См. также \citet[лемма~A.1]{rudomanov2026everest}
и \citet[лемма~2.5]{vankov2025optimizing}.

\begin{lemma}[рост градиента вдоль отрезка]\label{lem:local-transfer}
Пусть $\varphi\in C^2(\mathbb R^d)$ и
$\norm{\nabla^2\varphi(x)}\le L_0+L_1\norm{\nabla\varphi(x)}$
для всех $x$. Тогда для любых $x,y$ выполняется
\begin{equation}\label{eq:local-transfer}
 L_0+L_1\norm{\nabla\varphi(x)}
 \le e^{L_1\norm{x-y}}
       \bigl(L_0+L_1\norm{\nabla\varphi(y)}\bigr).
\end{equation}
В частности, на шаре $B(y,r)$ градиент $\varphi$ липшицев с константой
$e^{L_1r}(L_0+L_1\norm{\nabla\varphi(y)})$.
\end{lemma}

\begin{proof}
По формуле Ньютона--Лейбница
\[
 \nabla\varphi(y+t(x-y))-\nabla\varphi(y)
 =\int_0^t\nabla^2\varphi(y+u(x-y))(x-y)\,du.
\]
Взяв норму и применив условие на гессиан, для $t\in[0,1]$ получаем
\[
 L_0+L_1\norm{\nabla\varphi(y+t(x-y))}
 \le L_0+L_1\norm{\nabla\varphi(y)}
 +L_1\norm{x-y}\int_0^t\bigl(L_0+L_1\norm{\nabla\varphi(y+u(x-y))}\bigr)\,du.
\]
По неравенству Гронуолла
\[
 L_0+L_1\norm{\nabla\varphi(y+t(x-y))}
 \le e^{L_1\norm{x-y}t}\bigl(L_0+L_1\norm{\nabla\varphi(y)}\bigr).
\]
При $t=1$ получаем первое утверждение. Для $x\in B(y,r)$ оно даёт оценку
нормы гессиана на всём шаре, а интегрирование гессиана по отрезкам
внутри выпуклого шара доказывает липшицевость градиента.
\end{proof}

\citep[Теорема Байона--Хаддада,][]{bauschke2010}

\begin{lemma}[кокоэрцитивность на выпуклой области]\label{lem:local-cocoercivity}
Пусть $\varphi$ выпукла и дважды непрерывно дифференцируема
в окрестности выпуклого множества $\mathcal D$. Если
$\norm{\nabla^2\varphi}\le L$ на $\mathcal D$, где $L>0$, то для всех
$x,z\in\mathcal D$
\[
 \norm{\nabla\varphi(x)-\nabla\varphi(z)}^2
 \le L\ip{\nabla\varphi(x)-\nabla\varphi(z)}{x-z}.
\]
\end{lemma}

\begin{proof}
Средний гессиан
\[
 H=\int_0^1\nabla^2\varphi(z+t(x-z))\,dt
\]
удовлетворяет $0\preceq H\preceq LI$ по выпуклости, поэтому
$H^2\preceq LH$ и
\[
 \norm{\nabla\varphi(x)-\nabla\varphi(z)}^2
 =\ip{x-z}{H^2(x-z)}
 \le L\ip{x-z}{H(x-z)}
 =L\ip{\nabla\varphi(x)-\nabla\varphi(z)}{x-z}.
\]
\end{proof}

\begin{lemma}[оценка разности градиентов]\label{lem:local-bregman}
Пусть $\varphi$ выпукла, $\varphi\in C^2(\mathbb R^d)$ и
$\norm{\nabla^2\varphi}\le L$ на $B(y,3r)$, где $L>0$.
Тогда для $x,w\in B(y,r)$ выполнено
\[
 \norm{\nabla\varphi(x)-\nabla\varphi(w)}^2
 \le2L\bigl(\varphi(x)-\varphi(w)
                     -\ip{\nabla\varphi(w)}{x-w}\bigr).
\]
\end{lemma}

\begin{proof}
Так как градиент функции $\varphi$ является $L$-липшицевым на $B(y,r)$
и $x,w\in B(y,r)$, имеем
\[
 \norm{\nabla\varphi(x)-\nabla\varphi(w)}\le L\norm{x-w}\le2Lr.
\]
Поэтому точка
\[
 q=x-\frac1L\bigl(\nabla\varphi(x)-\nabla\varphi(w)\bigr)
\]
лежит в $B(y,3r)$, а значит, и весь отрезок $[x,q]$ содержится в $B(y,3r)$.
Функция
\[
 z\mapsto\varphi(z)-\varphi(w)-\ip{\nabla\varphi(w)}{z-w}
\]
выпукла, неотрицательна во всём пространстве и имеет на $B(y,3r)$
$L$-липшицев градиент $\nabla\varphi(z)-\nabla\varphi(w)$.
Квадратичная оценка сверху на отрезке $[x,q]$ даёт
\begin{align*}
 0&\le\varphi(q)-\varphi(w)-\ip{\nabla\varphi(w)}{q-w}
 \\&\le\varphi(x)-\varphi(w)-\ip{\nabla\varphi(w)}{x-w}
   +\ip{\nabla\varphi(x)-\nabla\varphi(w)}{q-x}+\frac L2\norm{q-x}^2
 \\&=\varphi(x)-\varphi(w)-\ip{\nabla\varphi(w)}{x-w}
   -\frac1{2L}\norm{\nabla\varphi(x)-\nabla\varphi(w)}^2.
\end{align*}
Остаётся перенести последнее слагаемое в левую часть.
\end{proof}

\subsection{Оценка нормы градиента через невязку}

Ср. \citet{li2023generalized} и \citet[лемма~3.4]{yu2026mirror}.

\begin{lemma}[интегральная оценка нормы градиента]\label{lem:self-bound}
Пусть $\varphi$ выпукла, ограничена снизу, принадлежит
$C^2(\mathbb R^d)$ и удовлетворяет условию леммы~\ref{lem:local-transfer}.
Тогда для любого $x$
\begin{equation}\label{eq:integral-self-bound}
 \begin{gathered}
 \int_0^{\norm{\nabla\varphi(x)}}\frac{t}{L_0+L_1t}\,dt
 \le\varphi(x)-\inf\varphi,\\
 \norm{\nabla\varphi(x)}
 \le L_1\bigl(\varphi(x)-\inf\varphi\bigr)
 +\sqrt{2L_0\bigl(\varphi(x)-\inf\varphi\bigr)}.
 \end{gathered}
\end{equation}
При $L_0=0$ подынтегральная функция доопределяется в нуле по непрерывности.
Достижимость инфимума не предполагается.
\end{lemma}

\begin{proof}
При $\nabla\varphi(x)=0$ утверждение очевидно. Иначе выберем $h>0$ так, чтобы
\[
 \begin{gathered}
 L_1h\norm{\nabla\varphi(x)}\le\tfrac12,\\
 h\bigl(L_0+L_1\norm{\nabla\varphi(x)}\bigr)\le\tfrac12,
 \end{gathered}
\]
и рассмотрим градиентный спуск из точки $x_0=x$:
\[
 x_{j+1}=x_j-h\nabla\varphi(x_j).
\]

\begin{itemize}
\item Норма градиента не возрастает. Пусть
$\norm{\nabla\varphi(x_j)}\le\norm{\nabla\varphi(x)}$. Отрезок
$[x_j,x_{j+1}]$ имеет длину $h\norm{\nabla\varphi(x_j)}$, поэтому
по лемме~\ref{lem:local-transfer} норма гессиана на нём не превосходит
\begin{equation}\label{eq:self-bound-step-curvature}
 e^{L_1h\norm{\nabla\varphi(x_j)}}
 \bigl(L_0+L_1\norm{\nabla\varphi(x_j)}\bigr)
 \le2\bigl(L_0+L_1\norm{\nabla\varphi(x)}\bigr)
 \le\frac1h.
\end{equation}
Значит, средний гессиан $H_j$ на этом отрезке удовлетворяет
$0\preceq hH_j\preceq I$, и
\[
 \begin{gathered}
 \nabla\varphi(x_{j+1})=(I-hH_j)\nabla\varphi(x_j),\\
 \norm{\nabla\varphi(x_{j+1})}\by{\le}{(a)}\norm{\nabla\varphi(x_j)}.
 \end{gathered}
\]
Здесь (a) следует из $0\preceq I-hH_j\preceq I$.
По индукции эти оценки верны для всех $j$.

\item Оценим две телескопические суммы. Квадратичная оценка сверху
вместе с оценкой кривизны~\eqref{eq:self-bound-step-curvature} даёт
\[
 \varphi(x_j)-\varphi(x_{j+1}) \ge h\norm{\nabla\varphi(x_j)}^2
 -\frac{h^2}{2}e^{L_1h\norm{\nabla\varphi(x_j)}}
 \bigl(L_0+L_1\norm{\nabla\varphi(x_j)}\bigr)\norm{\nabla\varphi(x_j)}^2 \ge\frac h2\norm{\nabla\varphi(x_j)}^2.
\]
Подынтегральная функция $t/(L_0+L_1t)$ не убывает, а
\[
 \norm{\nabla\varphi(x_j)}-\norm{\nabla\varphi(x_{j+1})}
 \le\norm{hH_j\nabla\varphi(x_j)}.
\]
Поэтому
\[
 \int_{\norm{\nabla\varphi(x_{j+1})}}^{\norm{\nabla\varphi(x_j)}}
 \frac{t\,dt}{L_0+L_1t} \le\frac{\norm{\nabla\varphi(x_j)}}{L_0+L_1\norm{\nabla\varphi(x_j)}}
 \norm{hH_j\nabla\varphi(x_j)} \by{\le}{(a)}
 he^{L_1h\norm{\nabla\varphi(x_j)}}\norm{\nabla\varphi(x_j)}^2.
\]
Здесь (a) следует из оценки~\eqref{eq:self-bound-step-curvature}
для $\norm{H_j}$. Так как $e^s-1\le2s$ и $e^s\le2$ при $0\le s\le1/2$,
из оценки убывания функции и последней оценки интеграла получаем
\[
 \int_{\norm{\nabla\varphi(x_{j+1})}}^{\norm{\nabla\varphi(x_j)}}
 \frac{t\,dt}{L_0+L_1t}
 \le\varphi(x_j)-\varphi(x_{j+1})
 +h^2\bigl(L_0+3L_1\norm{\nabla\varphi(x)}\bigr)\norm{\nabla\varphi(x_j)}^2.
\]
Суммируя неравенство
\[
 \varphi(x_i)-\varphi(x_{i+1})\ge\frac h2\norm{\nabla\varphi(x_i)}^2
\]
по $i=0,\ldots,m-1$, получаем
\[
 h\sum_{i<m}\norm{\nabla\varphi(x_i)}^2\le2\bigl(\varphi(x)-\inf\varphi\bigr).
\]
Суммируя оценку разности интегралов по $i=0,\ldots,j-1$, где $0\le j<m$,
и используя полученную оценку суммы квадратов норм градиентов, получаем
\[
 \int_0^{\norm{\nabla\varphi(x)}}\frac{t\,dt}{L_0+L_1t}
 \le\bigl(\varphi(x)-\inf\varphi\bigr)
 \Bigl(1+2h\bigl(L_0+3L_1\norm{\nabla\varphi(x)}\bigr)\Bigr)
 +\int_0^{\norm{\nabla\varphi(x_j)}}\frac{t\,dt}{L_0+L_1t}.
\]
Возьмём $m=\lceil h^{-2}\rceil$ и индекс $j<m$ с наименьшей нормой
$\norm{\nabla\varphi(x_j)}$. Тогда
\[
 \norm{\nabla\varphi(x_j)}^2
 \le\frac{2(\varphi(x)-\inf\varphi)}{mh}\to0
 \quad\text{при }h\downarrow0.
\]
Последний интеграл не превосходит $\norm{\nabla\varphi(x_j)}/L_1$,
поэтому, переходя к пределу, получаем первое
неравенство~\eqref{eq:integral-self-bound}.

\item При $L_0=0$ интеграл равен $\norm{\nabla\varphi(x)}/L_1$, что сразу
даёт второе неравенство. При $L_0>0$ после замены $t=L_0s/L_1$
интегральная оценка принимает вид
\[
 \frac{L_1\norm{\nabla\varphi(x)}}{L_0}
 -\log\Bigl(1+\frac{L_1\norm{\nabla\varphi(x)}}{L_0}\Bigr)
 \le\frac{L_1^2(\varphi(x)-\inf\varphi)}{L_0}.
\]
Функция $s\mapsto s-\log(1+s)$ возрастает при $s\ge0$. Для
$v\ge0$ из разложения экспоненты в ряд следует $e^{\sqrt{2v}}\ge1+\sqrt{2v}+v$, поэтому
\[
 (v+\sqrt{2v})-\log(1+v+\sqrt{2v})\ge v.
\]
Применив это неравенство при $v=L_1^2(\varphi(x)-\inf\varphi)/L_0$, получаем
\[
 \frac{L_1\norm{\nabla\varphi(x)}}{L_0}
 \le\frac{L_1^2(\varphi(x)-\inf\varphi)}{L_0}
 +\sqrt{\frac{2L_1^2(\varphi(x)-\inf\varphi)}{L_0}},
\]
что и требовалось.
\end{itemize}
\end{proof}

\subsection{Оценка локальной кривизны суммы}\label{app:observed-curvature}

\begin{proof}[Доказательство леммы~\ref{lem:observed-curvature}\proofup{lem:observed-curvature}]
Для среднего неотрицательных величин $f_i(y)-\inf f_i$ имеем
\[
 \frac1n\sum_i\bigl(f_i(y)-\inf f_i\bigr)
 =f(y)-f^\ast+f^\ast-\frac1n\sum_i\inf f_i
 \by{\le}{\text{\ref{asm:heterogeneity}}}\delta+f(y)-f^\ast.
\]
По лемме~\ref{lem:self-bound}, применимой в силу
предположений~\ref{asm:components} и~\ref{asm:smoothness},
\begin{align*}
 &\frac1n\sum_i\norm{\nabla f_i(y)}\le\frac{L_1}n\sum_i\bigl(f_i(y)-\inf f_i\bigr)
 +\frac1n\sum_i\sqrt{2L_{0,i}\bigl(f_i(y)-\inf f_i\bigr)}
 \\&\qquad\by{\le}{(a)}
 L_1\bigl(\delta+f(y)-f^\ast\bigr)+\sqrt{2L_0\bigl(\delta+f(y)-f^\ast\bigr)},
 \\
 &\ell(y)\le L_0+L_1^2\bigl(\delta+f(y)-f^\ast\bigr)
 +L_1\sqrt{2L_0\bigl(\delta+f(y)-f^\ast\bigr)}
 \\&\qquad\by{\le}{(b)}
 \tfrac32L_0+2L_1^2\bigl(\delta+f(y)-f^\ast\bigr).
\end{align*}
В (a) использованы неравенство Коши--Буняковского--Шварца
и ограничение на среднее значение $L_{0,i}$ из предположения~\ref{asm:smoothness};
(b) следует из $2ab\le a^2+b^2$.
Неравенство треугольника даёт обе оценки для $\norm{\nabla f(y)}$.
Остальные утверждения получаются суммированием оценок
леммы~\ref{lem:local-transfer} с учётом предположения~\ref{asm:smoothness}.
\end{proof}

\Needspace{10\baselineskip}
\subsection{Вырожденный класс}\label{app:degenerate}

\begin{proposition}\label{prop:degenerate-class}
При предположениях~\ref{asm:components}--\ref{asm:minimizer}
и~\ref{asm:heterogeneity}, если $L_0=\delta=0$, то все компоненты постоянны.
Следовательно, минимаксная оракульная сложность такого класса равна нулю.
\end{proposition}

\begin{proof}
В силу предположения~\ref{asm:smoothness} из $L_0=0$ следует
$L_{0,i}=0$. В точке $x^\ast$ из предположения~\ref{asm:minimizer} имеем
\[
 0\le\frac1n\sum_i\bigl(f_i(x^\ast)-\inf f_i\bigr)
 =f^\ast-\frac1n\sum_i\inf f_i\by{\le}{\text{\ref{asm:heterogeneity}}}\delta=0.
\]
Значит, каждое слагаемое равно нулю, откуда $\nabla f_i(x^\ast)=0$.
Лемма~\ref{lem:local-transfer} при $L_{0,i}=0$ даёт
\[
 \norm{\nabla f_i(x)}\le e^{L_1\norm{x-x^\ast}}
 \norm{\nabla f_i(x^\ast)}=0
\]
для любого $x$. Следовательно, все компоненты постоянны, и ответ $0$
можно выдать без единого запроса к оракулу.
\end{proof}

